% GENERATED FILE. Edit canonical lessons, then run npm run generate:books.
% canonical-source-hash: fnv1a64:04f2bfa0440f95c4
% canonical-lessons: UR-C541-zaviya, UR-C541-halat, UR-C541-mashrut, UR-C541-naubat, UR-C541-bilfarz

\chapter{In That Case}
\label{ch:ur-in-that-case}
\hlchaptermodality{541}

\begin{chapteropening}
\textbf{By the end of this chapter:} \emph{I can say what we will do in each case.}

Complete the last lesson of chapter 541: I can say what we will do in each case.
\end{chapteropening}

\section[\texorpdfstring{\ur{زاویہ} (zāviyā)}{زاویہ (zāviyā)}]{\ur{زاویہ} (zāviyā) — an angle; a point of view}

\label{lesson:UR-C541-zaviya}



\begin{quote}
Before the new one: say the Urdu for damages (to be paid), then the Urdu for to be robbed.
\end{quote}

\subsection*{You'll want to know: \ur{زاویہ}}
\textbf{\ur{زاویہ}} — \emph{zāviyā} — \textquotedblleft{}an angle; a point of view\textquotedblright{}.

\begin{quote}
\ur{اگر} \ur{آپ} \ur{اس} \ur{زاویے} \ur{سے} \ur{سوچتے} \ur{تو} \ur{بات} \ur{سمجھ} \ur{جاتے۔}
\end{quote}

\begin{quote}
\emph{agar āp is zāviye se sochte to bāt samajh jāte.}
\end{quote}

\begin{quote}
If you looked at it from this angle, you would understand.
\end{quote}

Masculine.

\begin{cousinweb}
\textbf{\ur{زاویہ}} (\emph{zāviyā}) is usually traced, through Persian, to Arabic \emph{zāwiya}, \textquotedblleft{}a corner, an angle\textquotedblright{}.
\end{cousinweb}

\begin{grammarlens}[title={What you've built}]
One more word for saying what happens if.
\end{grammarlens}

\subsection*{Your turn}
\begin{itemize}
\raggedright
  \item \emph{Say it:} \emph{zāviyā}
  \item \emph{Say it:} \emph{zāviyā}, once more
  \item \emph{Say it:} \emph{agar āp is zāviye se sochte to bāt samajh jāte.} — If you looked at it from this angle, you would understand.
  \item \emph{Recall:} say the Urdu for damages (to be paid), then the Urdu for to be robbed, then say \emph{zāviyā} again
\end{itemize}

\begin{tcolorbox}[breakable,colback=teal!4,colframe=teal!35!black,title={Before you move on}]
What does \ur{زاویہ} mean? (\textquotedblleft{}An angle; a point of view\textquotedblright{}.) Say it once more.
\end{tcolorbox}
\section[\texorpdfstring{\ur{حالات} (hālāt)}{حالات (hālāt)}]{\ur{حالات} (hālāt) — circumstances}

\label{lesson:UR-C541-halat}



\begin{quote}
Before the new one: say the Urdu for to be robbed, then the Urdu for an angle.
\end{quote}

\subsection*{You'll want to know: \ur{حالات}}
\textbf{\ur{حالات}} — \emph{hālāt} — \textquotedblleft{}circumstances\textquotedblright{}.

\begin{quote}
\ur{اگر} \ur{حالات} \ur{بدلے} \ur{تو} \ur{ہم} \ur{اپنا} \ur{پروگرام} \ur{بدلیں} \ur{گے۔}
\end{quote}

\begin{quote}
\emph{agar hālāt badle to ham apnā progrām badleṅ ge.}
\end{quote}

\begin{quote}
If circumstances change, we will change our plans.
\end{quote}

Masculine, used in the plural.

\begin{cousinweb}
\textbf{\ur{حالات}} (\emph{hālāt}) is usually traced, through Persian, to Arabic \emph{ḥālāt}, the plural of \emph{ḥāl}, \textquotedblleft{}a state, a condition\textquotedblright{}.
\end{cousinweb}

\begin{grammarlens}[title={What you've built}]
One more word for saying what happens if.
\end{grammarlens}

\subsection*{Your turn}
\begin{itemize}
\raggedright
  \item \emph{Say it:} \emph{hālāt}
  \item \emph{Say it:} \emph{hālāt}, once more
  \item \emph{Say it:} \emph{agar hālāt badle to ham apnā progrām badleṅ ge.} — If circumstances change, we will change our plans.
  \item \emph{Recall:} say the Urdu for to be robbed, then the Urdu for an angle, then say \emph{hālāt} again
\end{itemize}

\begin{tcolorbox}[breakable,colback=teal!4,colframe=teal!35!black,title={Before you move on}]
What does \ur{حالات} mean? (\textquotedblleft{}Circumstances\textquotedblright{}.) Say it once more.
\end{tcolorbox}
\section[\texorpdfstring{\ur{مشروط} (mashrūt)}{مشروط (mashrūt)}]{\ur{مشروط} (mashrūt) — conditional; with strings attached}

\label{lesson:UR-C541-mashrut}



\begin{quote}
Before the new one: say the Urdu for an angle, then the Urdu for circumstances.
\end{quote}

\subsection*{You'll want to know: \ur{مشروط}}
\textbf{\ur{مشروط}} — \emph{mashrūt} — \textquotedblleft{}conditional; with strings attached\textquotedblright{}.

\begin{quote}
\ur{اگر} \ur{ان} \ur{کی} \ur{ہاں} \ur{مشروط} \ur{ہے} \ur{تو} \ur{شرط} \ur{بتائیے۔}
\end{quote}

\begin{quote}
\emph{agar un kī hā̃ mashrūt hai to shart batāiye.}
\end{quote}

\begin{quote}
If their yes comes with a condition, tell me the condition.
\end{quote}

\begin{cousinweb}
\textbf{\ur{مشروط}} (\emph{mashrūt}) is usually traced, through Persian, to Arabic \emph{mashrūṭ}, \textquotedblleft{}made conditional\textquotedblright{}.
\end{cousinweb}

\begin{grammarlens}[title={What you've built}]
One more word for saying what happens if.
\end{grammarlens}

\subsection*{Your turn}
\begin{itemize}
\raggedright
  \item \emph{Say it:} \emph{mashrūt}
  \item \emph{Say it:} \emph{mashrūt}, once more
  \item \emph{Say it:} \emph{agar un kī hā̃ mashrūt hai to shart batāiye.} — If their yes comes with a condition, tell me the condition.
  \item \emph{Recall:} say the Urdu for an angle, then the Urdu for circumstances, then say \emph{mashrūt} again
\end{itemize}

\begin{tcolorbox}[breakable,colback=teal!4,colframe=teal!35!black,title={Before you move on}]
What does \ur{مشروط} mean? (\textquotedblleft{}Conditional; with strings attached\textquotedblright{}.) Say it once more.
\end{tcolorbox}
\section[\texorpdfstring{\ur{نوبت} (naubat)}{نوبت (naubat)}]{\ur{نوبت} (naubat) — a point (that things come to)}

\label{lesson:UR-C541-naubat}



\begin{quote}
Before the new one: say the Urdu for circumstances, then the Urdu for conditional.
\end{quote}

\subsection*{You'll want to know: \ur{نوبت}}
\textbf{\ur{نوبت}} — \emph{naubat} — \textquotedblleft{}a point (that things come to)\textquotedblright{}.

\begin{quote}
\ur{اگر} \ur{نوبت} \ur{یہاں} \ur{تک} \ur{پہنچی} \ur{تو} \ur{پولیس} \ur{کو} \ur{بتائیں} \ur{گے۔}
\end{quote}

\begin{quote}
\emph{agar naubat yahāṅ tak pahuṅchī to polīs ko batāeṅ ge.}
\end{quote}

\begin{quote}
If it comes to that, we will tell the police.
\end{quote}

Feminine.

\begin{cousinweb}
\textbf{\ur{نوبت}} (\emph{naubat}) is usually traced, through Persian, to Arabic \emph{nawba}, \textquotedblleft{}a turn, a time\textquotedblright{}.
\end{cousinweb}

\begin{grammarlens}[title={What you've built}]
One more word for saying what happens if.
\end{grammarlens}

\subsection*{Your turn}
\begin{itemize}
\raggedright
  \item \emph{Say it:} \emph{naubat}
  \item \emph{Say it:} \emph{naubat}, once more
  \item \emph{Say it:} \emph{agar naubat yahāṅ tak pahuṅchī to polīs ko batāeṅ ge.} — If it comes to that, we will tell the police.
  \item \emph{Recall:} say the Urdu for circumstances, then the Urdu for conditional, then say \emph{naubat} again
\end{itemize}

\begin{tcolorbox}[breakable,colback=teal!4,colframe=teal!35!black,title={Before you move on}]
What does \ur{نوبت} mean? (\textquotedblleft{}A point (that things come to)\textquotedblright{}.) Say it once more.
\end{tcolorbox}
\section[\texorpdfstring{\ur{بالفرض} (bilfarz)}{بالفرض (bilfarz)}]{\ur{بالفرض} (bilfarz) — supposing (for argument's sake)}

\label{lesson:UR-C541-bilfarz}



\begin{quote}
Before the new one: say the Urdu for conditional, then the Urdu for a point (that things come to).
\end{quote}

\subsection*{You'll want to know: \ur{بالفرض}}
\textbf{\ur{بالفرض}} — \emph{bilfarz} — \textquotedblleft{}supposing (for argument's sake)\textquotedblright{}.

\begin{quote}
\ur{بالفرض} \ur{اگر} \ur{بارش} \ur{ہوئی} \ur{تو} \ur{ہم} \ur{کیا} \ur{کریں} \ur{گے؟}
\end{quote}

\begin{quote}
\emph{bilfarz agar bārish huī to ham kyā kareṅ ge?}
\end{quote}

\begin{quote}
Supposing it rains, what will we do?
\end{quote}

\begin{cousinweb}
\textbf{\ur{بالفرض}} (\emph{bilfarz}) is usually traced, through Persian, to Arabic \emph{bi-l-farḍ}, \textquotedblleft{}by way of supposing\textquotedblright{}, built from \emph{bi}, \textquotedblleft{}by\textquotedblright{}, \emph{al}, \textquotedblleft{}the\textquotedblright{}, and \emph{farḍ}, the word Urdu also has as \textbf{\ur{فرض}} (\emph{farz}), \textquotedblleft{}a duty\textquotedblright{}.
\end{cousinweb}

\begin{grammarlens}[title={What you've built}]
One more word for saying what happens if.
\end{grammarlens}

\subsection*{Your turn}
\begin{itemize}
\raggedright
  \item \emph{Say it:} \emph{bilfarz}
  \item \emph{Say it:} \emph{bilfarz}, once more
  \item \emph{Say it:} \emph{bilfarz agar bārish huī to ham kyā kareṅ ge?} — Supposing it rains, what will we do?
  \item \emph{Recall:} say the Urdu for conditional, then the Urdu for a point (that things come to), then say \emph{bilfarz} again
\end{itemize}

\begin{tcolorbox}[breakable,colback=teal!4,colframe=teal!35!black,title={Before you move on}]
What does \ur{بالفرض} mean? (\textquotedblleft{}Supposing (for argument's sake)\textquotedblright{}.) Say it once more.
\end{tcolorbox}
